\documentclass[11pt, a4paper]{article}
\usepackage[a4paper, margin=2.2cm]{geometry}
\usepackage{fontspec}
\usepackage[vietnamese, english, bidi=basic, provide=*]{babel}
\babelprovide[import, onchar=ids fonts]{vietnamese}
\babelfont{rm}{Noto Sans}
\usepackage{amsmath, amssymb}
\usepackage{booktabs}
\usepackage{xcolor}
\usepackage{enumitem}
\usepackage{mdframed}
\usepackage{hyperref}

\definecolor{primary}{RGB}{30, 64, 175}
\definecolor{secondary}{RGB}{15, 118, 110}
\definecolor{bgcode}{RGB}{248, 250, 252}
\definecolor{bordercode}{RGB}{203, 213, 225}

\hypersetup{
    colorlinks=true,
    linkcolor=primary,
    urlcolor=primary
}

\mdfdefinestyle{codeblock}{
    backgroundcolor=bgcode,
    linecolor=bordercode,
    linewidth=1pt,
    topline=true,
    bottomline=true,
    leftline=true,
    rightline=true,
    innerleftmargin=10pt,
    innerrightmargin=10pt,
    innertopmargin=8pt,
    innerbottommargin=8pt
}

\title{\textbf{\Large BÁO CÁO CHI TIẾT QUESTION 4: NHẬN DIỆN NGÔN NGỮ BẰNG RNN}\\[0.3em]
\large Project 5: Machine Learning --- CS 188 (Fall 2025)}
\author{\textbf{Trí tuệ nhân tạo | UET}}
\date{\today}

\begin{document}
\maketitle

\section{Tổng quan về Question 4}
Trong \textbf{Project 5: Machine Learning} của khóa học CS 188 (Fall 2025), \textbf{Question 4} yêu cầu người học cài đặt mô hình \textbf{LanguageIDModel} trong tệp \texttt{models.py}. Bài toán đặt ra là \textbf{nhận diện ngôn ngữ của một từ} chỉ dựa trên chuỗi ký tự tạo nên từ đó. Vì các từ có độ dài khác nhau, mô hình cần xử lý dữ liệu dạng chuỗi, và đây là lý do \textbf{mạng nơ-ron hồi quy (Recurrent Neural Network -- RNN)} được sử dụng.

\section{Kiến thức sử dụng (What Knowledge is Used?)}
Để thực hiện thành công Question 4, người học cần áp dụng các khối kiến thức sau:
\begin{itemize}[itemsep=4pt]
    \item \textbf{Học có giám sát (Supervised Learning):}
    \begin{itemize}
        \item Mô hình học từ tập dữ liệu gồm các từ đã được gán nhãn ngôn ngữ tương ứng.
        \item Bài toán thuộc dạng phân loại đa lớp (multi-class classification).
    \end{itemize}
    \item \textbf{Mạng nơ-ron hồi quy (Recurrent Neural Network -- RNN):}
    \begin{itemize}
        \item Xử lý tuần tự từng ký tự của chuỗi đầu vào.
        \item Duy trì \textbf{trạng thái ẩn (hidden state)} để lưu giữ thông tin từ những ký tự đã xuất hiện trước đó.
    \end{itemize}
    \item \textbf{Biểu diễn dữ liệu (Data Representation):}
    \begin{itemize}
        \item Chuyển các ký tự thành dạng số để mô hình có thể xử lý, ví dụ bằng \textbf{One-hot Encoding}.
    \end{itemize}
    \item \textbf{Hàm mất mát Cross-Entropy:}
    \begin{itemize}
        \item Đo lường sự khác biệt giữa dự đoán của mô hình và nhãn ngôn ngữ thực tế.
    \end{itemize}
    \item \textbf{Lan truyền ngược qua thời gian (Backpropagation Through Time -- BPTT):}
    \begin{itemize}
        \item Lan truyền ngược gradient qua các bước thời gian để cập nhật trọng số của RNN trong quá trình huấn luyện.
    \end{itemize}
\end{itemize}

\section{Bài toán giải quyết \& Ứng dụng thực tế}

\subsection{Bài toán giải quyết}
Question 4 giải quyết bài toán \textbf{Nhận diện ngôn ngữ (Language Identification)}: Cho trước một từ được biểu diễn bởi chuỗi ký tự $(c_1, c_2, \dots, c_T)$, mục tiêu là xác định từ đó thuộc ngôn ngữ nào trong tập các ngôn ngữ cho trước. Do độ dài chuỗi $T$ thay đổi theo từng từ, mô hình phải đọc lần lượt từng ký tự và tổng hợp thông tin của toàn bộ chuỗi trước khi đưa ra kết luận.

\subsection{Ứng dụng thực tế}
Tư duy nhận diện ngôn ngữ từ chuỗi ký tự được áp dụng rộng rãi trong thực tiễn:
\begin{itemize}[itemsep=3pt]
    \item \textbf{Nhận diện ngôn ngữ từ từ, văn bản (Language Detection):} Tự động xác định ngôn ngữ của một từ hoặc đoạn văn bản bất kỳ.
    \item \textbf{Phân loại văn bản đa ngôn ngữ (Multilingual Text Classification):} Phân nhóm tài liệu theo ngôn ngữ trước khi xử lý hoặc lưu trữ.
    \item \textbf{Hỗ trợ hệ thống xử lý ngôn ngữ tự nhiên (NLP) và tiền xử lý dữ liệu văn bản:} Lựa chọn đúng bộ công cụ (dịch máy, kiểm tra chính tả, phân tích cú pháp) cho từng ngôn ngữ.
\end{itemize}

\section{Công thức toán học \& Quy trình Cài đặt (Implementation)}

\subsection{1. Biểu diễn đầu vào (Input Representation)}
Mỗi ký tự $c_t$ được chuyển thành một vector số $x_t$ bằng \textbf{One-hot Encoding}: vector có độ dài bằng kích thước bảng ký tự, trong đó chỉ phần tử ứng với ký tự đó có giá trị $1$, các phần tử còn lại bằng $0$. Chuỗi vector $(x_1, x_2, \dots, x_T)$ sau đó được đưa lần lượt vào RNN.

\subsection{2. Cập nhật trạng thái ẩn (Hidden State Update)}
Tại mỗi bước thời gian $t$, RNN kết hợp thông tin từ ký tự hiện tại $x_t$ với trạng thái ẩn trước đó $h_{t-1}$:
\begin{equation}
h_t = \text{ReLU}\left(x_t W_x + h_{t-1} W_h + b_h\right)
\end{equation}
Trong đó:
\begin{itemize}[itemsep=2pt]
    \item $x_t$: vector biểu diễn ký tự tại vị trí $t$.
    \item $h_{t-1}$: trạng thái ẩn của bước trước (với $h_0$ khởi tạo bằng vector không).
    \item $W_x, W_h$: các ma trận trọng số học được.
    \item $b_h$: vector bias.
    \item $h_t$: trạng thái ẩn mới, chứa thông tin được tích lũy trong quá trình xử lý chuỗi.
\end{itemize}

\subsection{3. Dự đoán ngôn ngữ (Language Prediction)}
Sau khi xử lý hết chuỗi ký tự, trạng thái ẩn cuối cùng $h_T$ được dùng để tạo ra điểm số (logits) cho từng lớp ngôn ngữ:
\begin{equation}
z = h_T W_o + b_o
\end{equation}
Trong đó $z$ là vector điểm số dự đoán, $W_o$ và $b_o$ lần lượt là ma trận trọng số và vector bias của tầng đầu ra. Mô hình chọn ngôn ngữ có điểm số cao nhất:
\begin{equation}
\hat{y} = \underset{i}{\arg\max}\; z_i
\end{equation}

\subsection{4. Huấn luyện mô hình (Training)}
Mô hình so sánh dự đoán với nhãn thực tế bằng hàm mất mát \textbf{Cross-Entropy}:
\begin{equation}
L = -\log P\left(y_{\text{true}} \mid x\right), \qquad P(i \mid x) = \frac{e^{z_i}}{\sum_j e^{z_j}}
\end{equation}
Sau đó, \textbf{lan truyền ngược qua thời gian (BPTT)} được sử dụng để tính gradient của $L$ theo các tham số $W_x, W_h, b_h, W_o, b_o$ và cập nhật trọng số nhằm giảm giá trị hàm mất mát.

\subsection{5. Cấu trúc hàm cần cài đặt trong \texttt{models.py}}
Trong bài tập CS 188, người học cần hoàn thiện lớp \texttt{LanguageIDModel} với các phương thức chính:
\begin{enumerate}[itemsep=3pt]
    \item \texttt{\_\_init\_\_(self)}: Khởi tạo các tham số của RNN ($W_x$, $W_h$, $b_h$) và tầng đầu ra ($W_o$, $b_o$).
    \item \texttt{run(self, xs)}: Nhận một chuỗi các tensor ký tự \texttt{xs}, cập nhật trạng thái ẩn qua từng bước thời gian và trả về tensor điểm số $z$ cho các lớp ngôn ngữ.
    \item \texttt{get\_loss(self, xs, y)}: Tính hàm mất mát Cross-Entropy giữa điểm số dự đoán và nhãn thực tế \texttt{y}.
    \item \texttt{train(self, dataset)}: Huấn luyện mô hình bằng gradient descent (qua BPTT) cho đến khi đạt độ chính xác yêu cầu trên tập kiểm tra.
\end{enumerate}

\section{Thống nhất Output \& Mô tả Bài toán (Problem Statement)}

\begin{mdframed}[style=codeblock]
\textbf{MÔ TẢ BÀI TOÁN CHUẨN ĐỂ ĐƯA VÀO BÁO CÁO / ĐỀ BÀI}

\textbf{Tên bài toán:} Nhận diện ngôn ngữ của một từ dựa trên chuỗi ký tự bằng Mạng nơ-ron hồi quy (\texttt{LanguageIDModel}).

\textbf{Đầu vào (Inputs):}
\begin{itemize}[itemsep=2pt]
    \item Tập dữ liệu huấn luyện \texttt{dataset} chứa các cặp $(\texttt{xs}, y)$.
    \item \texttt{xs}: chuỗi các tensor one-hot, mỗi tensor biểu diễn một ký tự của từ theo thứ tự xuất hiện.
    \item $y$: nhãn ngôn ngữ thực tế của từ, được mã hóa one-hot trên tập các ngôn ngữ.
\end{itemize}

\textbf{Đầu ra (Outputs):}
\begin{itemize}[itemsep=2pt]
    \item Phương thức \texttt{run(self, xs)}: Trả về tensor điểm số $z$, mỗi phần tử ứng với một lớp ngôn ngữ.
    \item Phương thức \texttt{get\_loss(self, xs, y)}: Trả về giá trị mất mát Cross-Entropy $L = -\log P(y_{\text{true}} \mid x)$.
    \item Phương thức \texttt{train(self, dataset)}: Một mô hình RNN đã được huấn luyện hoàn chỉnh, dự đoán ngôn ngữ tương ứng với từ đầu vào bằng cách chọn lớp có điểm số cao nhất $\hat{y} = \arg\max_i z_i$.
\end{itemize}
\end{mdframed}

\section{Kết luận}
Question 4 giúp người học làm quen với việc xử lý dữ liệu dạng chuỗi: biểu diễn ký tự bằng one-hot, duy trì thông tin qua trạng thái ẩn của RNN, và huấn luyện bằng Cross-Entropy cùng BPTT. Đây là nền tảng quan trọng để tiếp cận các mô hình xử lý ngôn ngữ tự nhiên phức tạp hơn như LSTM, GRU hay Transformer.

\end{document}